
\makeatletter
\let\@oldaddcontentsline\addcontentsline
\renewcommand{\addcontentsline}[3]{}
\section*{附录}
\let\addcontentsline\@oldaddcontentsline
\makeatother

\begin{table}[H]
  \centering
  \caption{支撑材料文件列表}
  \label{tab:appendix}
  \renewcommand{\arraystretch}{1.2}
  \begin{tabularx}{\textwidth}{LX}
    \toprule
    文件 & 说明 \\
    \midrule
    问题一\_xxx.py & 问题1求解代码 \\
    xxx.png & xxx图 \\
    \bottomrule
  \end{tabularx}
\end{table}

% ============================================================
% 源程序代码（2026规范第五条：附录应包括建模所用到的全部完整、可运行的
% 源程序代码，含EXCEL、SPSS等软件的交互命令；若确实没有用到程序，请将
% 本节替换为“本论文没有用到程序”；若没有支撑材料，请注明“本论文没有支撑材料”）
% 引入代码文件可用：\lstinputlisting[language=Python]{代码文件路径}
\subsection*{源程序代码}

\begin{lstlisting}[language=Python]
# 请在此处粘贴完整、可运行的源程序代码
\end{lstlisting}
